\documentclass[11pt]{article}

\usepackage{amsmath}
\usepackage[margin=1in]{geometry}

\title{Why Use Half Mean Squared Error?}
\author{CPE 342 Machine Learning}
\date{}

\begin{document}

\maketitle

\section{Definition}

Let the residual for observation $i$ be
\[
r_i = \hat{y}_i - y_i.
\]
The standard mean squared error (MSE) is
\[
J(c_0,c_1,c_2) = \frac{1}{n}\sum_{i=1}^{n} r_i^2.
\]
In this assignment, we use the half-MSE loss
\[
L(c_0,c_1,c_2) = \frac{1}{2n}\sum_{i=1}^{n} r_i^2
                 = \frac{1}{2}J(c_0,c_1,c_2).
\]

The factor $1/2$ does not change the fitted model. It only rescales the loss
and its gradient by a constant.

\section{Why the factor $1/2$ is useful}

The derivative of a squared residual contains a factor of $2$:
\[
\frac{\partial}{\partial \theta}r_i^2
= 2r_i\frac{\partial r_i}{\partial \theta}.
\]
For the half-MSE loss, the factor $2$ cancels the factor $1/2$:
\[
\frac{\partial L}{\partial \theta}
= \frac{1}{n}\sum_{i=1}^{n}
r_i\frac{\partial r_i}{\partial \theta}.
\]
This makes the gradient equations shorter and reduces the chance of carrying
an unnecessary factor of $2$ through every update rule.

\section{Gradients for the assignment model}

The model is
\[
\hat{y}_i = c_0 + c_1 e^{c_2x_i},
\]
so
\[
r_i = c_0 + c_1 e^{c_2x_i} - y_i.
\]
The half-MSE gradients are therefore
\begin{align*}
\frac{\partial L}{\partial c_0}
&= \frac{1}{n}\sum_{i=1}^{n}r_i, \\
\frac{\partial L}{\partial c_1}
&= \frac{1}{n}\sum_{i=1}^{n}r_i e^{c_2x_i}, \\
\frac{\partial L}{\partial c_2}
&= \frac{1}{n}\sum_{i=1}^{n}r_i c_1x_i e^{c_2x_i}.
\end{align*}

These gradients are the ones used by the notebook's batch gradient descent.

\section{Same optimum as standard MSE}

Because $L = J/2$ and $1/2$ is positive,
\[
\operatorname*{arg\,min}_{c_0,c_1,c_2} L
= \operatorname*{arg\,min}_{c_0,c_1,c_2} J.
\]
Therefore, both losses produce the same optimal coefficients, predictions,
residuals, and goodness-of-fit measures when optimized correctly.

The numerical loss values differ by a factor of two. For the notebook result,
the final half-MSE is approximately
\[
L = 0.0003777068,
\]
while the corresponding standard MSE is
\[
J = 2L = 0.0007554135.
\]
The normalization must be stated when reporting the loss value.

\section{Effect on the learning rate}

The gradients satisfy
\[
\nabla L = \frac{1}{2}\nabla J.
\]
Consequently, a half-MSE update with learning rate $\eta_L$ is
\[
\boldsymbol{c}^{(t+1)}
= \boldsymbol{c}^{(t)} - \eta_L\nabla L
= \boldsymbol{c}^{(t)} - \frac{\eta_L}{2}\nabla J.
\]
To obtain the same step as standard MSE gradient descent with learning rate
$\eta_J$, choose
\[
\eta_L = 2\eta_J.
\]
In the notebook, $\eta_L=0.2$ is used with half-MSE. The equivalent standard
MSE learning rate would be $\eta_J=0.1$.

\section{Conclusion}

Half-MSE is used for algebraic convenience. It removes the factor of $2$ from
the gradients without changing the optimization target or the final fitted
model. The only differences are the reported loss scale and the learning-rate
scale needed to produce equivalent gradient-descent steps.

\end{document}
